\hypertarget{aba-problem-memory-reclamation}{%
\chapter{ABA PROBLEM \& MEMORY
RECLAMATION}\label{aba-problem-memory-reclamation}}

In the rarefied air of lock-free programming, the \textbf{ABA Problem}
is the dragon that guards the gate. Conquering it requires understanding
the very fabric of memory lifecycles.

\hypertarget{the-aba-problem-explained}{%
\subsection{1. The ABA Problem
Explained}\label{the-aba-problem-explained}}

The Compare-And-Swap (CAS) primitive
(\passthrough{\lstinline!std::atomic::compare\_exchange\_weak!}) checks
if a value is \emph{equal} to an expected value. It does \textbf{not}
check if it is the \emph{same} object.

\textbf{The Scenario:} 1. Thread 1 reads pointer
\passthrough{\lstinline!A!} from a lock-free stack top. 2. Thread 1 is
preempted. 3. Thread 2 pops \passthrough{\lstinline!A!}, frees it, then
pushes \passthrough{\lstinline!B!}, then pushes a \emph{new} object
allocated at address \passthrough{\lstinline!A!} (recycled memory). 4.
Thread 1 wakes up, performs CAS. The address is still
\passthrough{\lstinline!A!}. CAS succeeds. 5. \textbf{Catastrophe:}
Thread 1 has popped the \emph{new} \passthrough{\lstinline!A!}, but its
local logic assumes it's the \emph{old} \passthrough{\lstinline!A!}
(e.g., pointing to \passthrough{\lstinline!B!} as the next node). The
stack is now corrupted.

\hypertarget{solution-i-tagged-pointers-version-counters}{%
\subsection{2. Solution I: Tagged Pointers (Version
Counters)}\label{solution-i-tagged-pointers-version-counters}}

Pack a version counter into the unused bits of a pointer (usually top 16
bits on 64-bit systems). * \textbf{Mechanism:} Every modification
increments the counter. \passthrough{\lstinline!Ptr(A, v1)!} !=
\passthrough{\lstinline!Ptr(A, v2)!}. * \textbf{Limitation:} Reduces
addressable memory space; requires platform-specific bit manipulation.

\begin{lstlisting}[language={C++}]
// Example of a 64-bit tagged pointer
struct TaggedPtr {
    uint64_t data; // 48 bits pointer, 16 bits tag

    TaggedPtr(void* ptr, uint16_t tag) {
        data = (reinterpret_cast<uint64_t>(ptr) & 0x0000FFFFFFFFFFFF) | (static_cast<uint64_t>(tag) << 48);
    }
    
    void* get_ptr() const { return reinterpret_cast<void*>(data & 0x0000FFFFFFFFFFFF); }
    uint16_t get_tag() const { return static_cast<uint16_t>(data >> 48); }
};
\end{lstlisting}

\hypertarget{solution-ii-hazard-pointers-the-gold-standard}{%
\subsection{3. Solution II: Hazard Pointers (The Gold
Standard)}\label{solution-ii-hazard-pointers-the-gold-standard}}

A \textbf{Hazard Pointer (HP)} is a thread-local signal saying ``I am
reading this object, do not delete it.''

\textbf{The Protocol:} 1. \textbf{Reader:} publish the pointer
\passthrough{\lstinline!P!} to a thread-local HP slot. 2.
\textbf{Reader:} Verify \passthrough{\lstinline!P!} is still in the data
structure. If not, retry. 3. \textbf{Writer (Deleter):} Unlink
\passthrough{\lstinline!P!} from the structure. 4. \textbf{Writer:}
Check all other threads' HPs. * If \passthrough{\lstinline!P!} is found
in any HP, add \passthrough{\lstinline!P!} to a ``Retire List'' (do not
\passthrough{\lstinline!delete!} yet). * If \passthrough{\lstinline!P!}
is not found, \passthrough{\lstinline!delete!} immediately. 5.
\textbf{Cleanup:} Periodically scan the Retire List and free objects no
longer protected by HPs.

\begin{itemize}
\tightlist
\item
  \textbf{Pros:} Wait-free readers, deterministic memory bound.
\item
  \textbf{Cons:} Heavy memory barrier usage (Store-Load fence needed
  after publishing HP).
\end{itemize}

\hypertarget{solution-iii-epoch-based-reclamation-ebr}{%
\subsection{4. Solution III: Epoch-Based Reclamation
(EBR)}\label{solution-iii-epoch-based-reclamation-ebr}}

Used by \passthrough{\lstinline!malloc!} implementations and databases
(like Silo). * \textbf{Concept:} A Global Epoch counter (E) and
per-thread Local Epochs (e\_t). * \textbf{Operation:} 1. Global Epoch
\passthrough{\lstinline!E!} increments periodically. 2. Threads update
\passthrough{\lstinline!e\_t = E!} when entering a critical section. 3.
Objects retired in Epoch \passthrough{\lstinline!E!} can be safely
deleted when all threads have reached Epoch
\passthrough{\lstinline!E+1!} or higher. * \textbf{Pros:} Extremely fast
(just checking integers). * \textbf{Cons:} One stalled thread prevents
\emph{all} memory reclamation (OOM risk).

\begin{center}\rule{0.5\linewidth}{0.5pt}\end{center}
